\clearpage\thispagestyle{empty}
\pdfbookmark[0]{Quelle und Übertragung}{connaissance1813-sources}
\begin{center}\textbf{Quelle und Übertragung}\end{center}
\hypertarget{connaissance1813-source-note}{}\label{connaissance1813-source-note}
\noindent KI-Übertragung des deutschen Originals von Carl Friedrich Gauß in modernes Deutsch. Die Besprechung der \emph{Connaissance des tems} für 1813 und 1814 ist auf den 3. Juni 1813 datiert. Die Übertragung folgt den \href{https://gdz.sub.uni-goettingen.de/id/PPN236007467}{\emph{Werken}, Band VI}, Göttingen 1874, S.\ 554--555. Der Erstdruck in den \emph{Göttingischen gelehrten Anzeigen}, Stück 88, S.\ 873--876, nennt Paris und Courcier vor den Titeln, gibt für 1814 das Oktavformat an und setzt nach \emph{bemerkten} ein Komma. Die Werkausgabe hat Stück 8 und einen Punkt nach \emph{bemerkten}. Die Übertragung bewahrt die Vorzeichen der Koeffizienten und die Wurzelnotation der Werkausgabe.

\noindent Das Bureau des Longitudes veröffentlichte beide Jahrgänge. Auf den Titelblättern der Exemplare der Bayerischen Staatsbibliothek stehen für \href{https://api.digitale-sammlungen.de/iiif/image/v2/bsb10538233_00005/full/1450,/0/default.jpg}{1813} Paris, Courcier, Juli 1811 und für \href{https://api.digitale-sammlungen.de/iiif/image/v2/bsb10538234_00005/full/1450,/0/default.jpg}{1814} Paris, Madame Veuve Courcier, April 1812. Der Jahrgang 1813 enthält Laplaces Aufsatz über die Kombination zahlreicher Beobachtungen (\href{https://api.digitale-sammlungen.de/iiif/image/v2/bsb10538233_00217/full/1450,/0/default.jpg}{S.\ 213}) und seinen Aufsatz über die langperiodische Ungleichheit des Mondlaufs (\href{https://api.digitale-sammlungen.de/iiif/image/v2/bsb10538233_00227/full/1450,/0/default.jpg}{S.\ 223}). Laplace druckt den Mondkoeffizienten als $-13''{,}92$ (\href{https://api.digitale-sammlungen.de/iiif/image/v2/bsb10538233_00229/full/1450,/0/default.jpg}{S.\ 225}). In Gauß' Besprechung steht $-13''96$. Lagranges Aufsatz über den Ursprung der Kometen im Jahrgang 1814 beginnt auf \href{https://api.digitale-sammlungen.de/iiif/image/v2/bsb10538234_00215/full/1450,/0/default.jpg}{S.\ 211}, unterscheidet auf \href{https://api.digitale-sammlungen.de/iiif/image/v2/bsb10538234_00217/full/1450,/0/default.jpg}{S.\ 213} rechtläufige und rückläufige Bahnen und nennt auf \href{https://api.digitale-sammlungen.de/iiif/image/v2/bsb10538234_00219/full/1450,/0/default.jpg}{S.\ 215} die Entfernung und die Geschwindigkeiten des Beispiels.

\noindent Erstdruck, digitalisiert von der Niedersächsischen Staats- und Universitätsbibliothek Göttingen, \href{https://images.p.sub.uni-goettingen.de/iiif/image/gdz:PPN319721507_1813:00000877/full/1650,/0/default.jpg}{S.\ 873}, \href{https://images.p.sub.uni-goettingen.de/iiif/image/gdz:PPN319721507_1813:00000878/full/1650,/0/default.jpg}{S.\ 874}, \href{https://images.p.sub.uni-goettingen.de/iiif/image/gdz:PPN319721507_1813:00000879/full/1650,/0/default.jpg}{S.\ 875}, \href{https://images.p.sub.uni-goettingen.de/iiif/image/gdz:PPN319721507_1813:00000880/full/1650,/0/default.jpg}{S.\ 876}.
\clearpage
